\SourceLeaf{335}

Sind Werte von $\mathfrak{B}, \mathfrak{C}$ vorgegeben, die diesen Bedingungen oder einer von ihnen nicht genügen, so steht fest, dass diese Aufgabe eine physische Unmöglichkeit impliziert.

Umgekehrt ist die Aufgabe möglich, wenn die aufgegebenen Werte den Bedingungen genügen, d. i. die beiden Hilfskreise treffen sich in einem Punkt, der für beide Kreise im geeigneten Abschnitt liegt. Sind nämlich 1) die Zeichen jener Sinus $+\;+\;+$, so kann jener Schnittpunkt nicht außerhalb des Hauptkreises liegen, wo die ersten beiden Sinus negativ, der dritte positiv sein müsste. Die beiden ersten Bedingungen würden nämlich erfordern, dass $\mathfrak{B}+180^0$, $\mathfrak{C}+180^0$ an die Stelle von $\mathfrak{B}, \mathfrak{C}$ treten, was aber wieder der dritten Bedingung widerspricht. Also liegt jener Schnittpunkt innerhalb des Hauptkreises, wo eben die geeigneten Abschnitte liegen.

Auf ähnliche Weise wird 2) für die Zeichen $-\;-\;+$ bewiesen, dass der Schnittpunkt nur außerhalb liegen kann, wo dann die geeigneten Abschnitte liegen.

\begin{center}
\begin{tabular}{c@{\hspace{2em}}cc}
Sind die Zeichen & \multicolumn{2}{c}{so liegt der Schnittpunkt} \\
 & im geeigneten Abschnitt & im ungeeigneten \\
$+\;+\;+$ & Kreis auf $a$ & [Kreis auf] $b, c$ \\
$+\;-\;+$ & $b$ & $a, c$ \\
$+\;-\;-$ & $c$ & $a, b$ \\
$-\;-\;-$ & $a$ & $b, c$ \\
$-\;+\;+$ & $c$ & $a, b$ \\
$-\;+\;-$ & $b$ & $a, c$ \\
\end{tabular}
\end{center}

Die elementare Konstruktion für das \textsc{Pothenot}sche Problem besteht darin, dass man Punkte in der durch (0) gehenden, zu 10 normalen Geraden bestimmt. Von einem solchen Punkt sind die Polarkoordinaten relativ zu (0):
\[ \text{Distanz } \frac{(12)}{\sin(02-01)}, \qquad \text{Richtung } 12+90^0-(02-01). \]

Eine andere neue Konstruktionsart besteht darin, dass man zunächst nicht den Punkt (0) sucht, sondern einen Punkt $(*)$, der in der Geraden 10 liegt und in der Distanz $(1*)=\frac{1}{(10)}$. Dieser Punkt wird durch Schnitte gerader Linien bestimmt, eine weniger, als von (0) aus Objekte geschnitten sind. Eine dieser
